\documentclass[11pt]{article}
\usepackage[a4paper,margin=25mm]{geometry}
\usepackage{amsmath,amssymb,amsthm,mathtools,hyperref}
\newcommand{\C}{\mathbb C}
\newcommand{\im}{\operatorname{im}}
\newcommand{\Tr}{\operatorname{Tr}}
\newcommand{\norm}[1]{\left\lVert#1\right\rVert}
\newtheorem{theorem}{Theorem}
\title{Adjacent original observation responses and their exact metric evolution}
\author{Split-Zero research programme}\date{19 September 2026}
\begin{document}\maketitle

\section{The unchanged packet, source and observation}
Fix the original actual simple quartet and its complete arithmetic unit,
period, logarithm branch and original observation. Keep
$a\equiv1\pmod4$, $q=(a+1)^2$, $c=a/2$ and the full polynomial
\[
 \chi(S)=\prod_{u,b=0}^a
 [S-c-(2u-a)\delta-i(2b-a)\gamma],\qquad
 0<\delta<1/2,\quad\gamma>2.                         \tag{ORE1}
\]
The same proof of the finite update applies to the complete primary
polynomial with multiplicities, retaining its full matrix dimension;
the source-specific references below retain their own stated domains.
Let $E=\C[S]/(\chi)$ in its original increasing-power frame.
The source norm is the literal integral of $|p(c+iy)|^2$ against
$m_a=w_h^{*a}(y)\,dy$, with $w_h=|(2\xi/h)(1/2+iy)|^2/(2\pi)$.
Let $G_N$ be the quotient metric induced from all source polynomials
of degree at most $N$, minimizing over all relations $\chi h$.
Let $p_n(S)$ be the original monic source orthogonal polynomial,
$\omega_n=\norm{p_n}_{m_a}^2>0$, and $b_n=[p_n]\in E$.
No Gamma source, rescaled mass or changed physical coordinate is substituted.

Retain the full original observation $\Lambda:E\to B=\im\Lambda$,
its full kernel $K$, and an injective frame $I:\C^{\dim K}\to E$.
The observation includes its original period map. In particular $I$
does not mean the native coefficient section $L_I$ of COI11.
Define, for $N\ge q-1$,
\begin{gather}
 H_N=I^*G_NI,\qquad
 P_N=I H_N^{-1}I^*G_N,\qquad
 Q_N=(\Lambda G_N^{-1}\Lambda^*)^{-1},\notag\\
 E_N=v^*G_Nv,\qquad
 \theta_N=\frac{v^*G_NP_Nv}{E_N},\qquad
 v=v_\lambda,\quad Av=\lambda v,\quad
 \lambda=c+a\delta+ia\gamma.                         \tag{ORE2}
\end{gather}
Here $A$ is original multiplication by $S$, $v$ is the full retained
class, and neither its low nor its conductor coordinate is discarded.
For $K=0$, $H_N$ is the empty matrix, its determinant is one and
$P_N=0$. For $K=E$, $B=0$ and the corresponding $Q_N$ is empty.
These conventions give the same formulas below.
The exact minimum-lift map is
$G_N^{-1}\Lambda^*Q_N:B\to E$, so
\[
 G_N(I-P_N)=\Lambda^*Q_N\Lambda,
 \qquad E_N(1-\theta_N)=(\Lambda v)^*Q_N(\Lambda v).    \tag{ORE3}
\]
This is the equality-constrained quadratic minimum described by
Boyd and Vandenberghe \cite[Appendix C.4.1, p.~673]{Boyd};
the following calculation supplies its complex Hermitian dictionary
for the actual $G_N$ and $\Lambda$.
Indeed the lift has observation equal to the prescribed value, is
orthogonal to $K$, and hence is its unique minimum. This proves both
identities on the original space, including its full cross metric.

\section{A single actual source increment and both determinant allocations}
At the following cutoff put $b=b_{N+1}$, $\omega=\omega_{N+1}$.
The exact source orthogonal decomposition
$\mathcal P_{N+1}=\mathcal P_N\mathbin{\perp}\C p_{N+1}$ implies
\begin{equation}\tag{ORE4}
 G_{N+1}^{-1}=G_N^{-1}+\frac{bb^*}{\omega},\qquad
 G_{N+1}=G_N-\frac{G_Nbb^*G_N}{\omega+\beta_N},
 \quad\beta_N=b^*G_Nb.
\end{equation}
For the first equality, the squared dual norm of a functional on $E$
is the sum of its squared values on a source orthogonal basis divided
by the original $\omega_n$; the new term is exactly $bb^*/\omega$.
For the second, multiply the proposed right side by the first matrix.
The coefficient of $G_Nbb^*/\omega$ cancels, proving the inverse
without any assumed metric equivalence. This is the exact rank-one
inverse identity recorded by Hager \cite[p.~221, equation (2)]{Hager},
with its source and denominator derived here.

On the original retained eigenclass domain NCE9 proves
$b^*G_Nv\ne0$, so $\beta_N>0$. Define
\begin{gather}
 r_{N+1}=\frac{\omega}{\omega+\beta_N}
        =\frac{\det G_{N+1}}{\det G_N},\qquad
 \xi_N=1-r_{N+1}\in(0,1),\notag\\
 \alpha_N=\frac{b^*G_NP_Nb}{\beta_N}\in[0,1],\qquad
 d_{K,N}=\frac{\det H_{N+1}}{\det H_N}.                \tag{ORE5}
\end{gather}
The determinant lemma follows by taking a basis whose first vector is
the single nonzero column of the update, or by expansion of the
determinant in that rank-one column. Applied to $H_N$ and (ORE4), it gives
\begin{equation}\tag{ORE6}
 \boxed{d_{K,N}=1-\xi_N\alpha_N,\qquad
 \frac{\det Q_{N+1}}{\det Q_N}=
       \frac{r_{N+1}}{d_{K,N}}.}
\end{equation}
For the second equation use (ORE3) with $b$:
$(\Lambda b)^*Q_N\Lambda b=\beta_N(1-\alpha_N)$.
Updating $Q_N^{-1}$ by $(\Lambda b)(\Lambda b)^*/\omega$
then proves the equation with denominator
$\omega+\beta_N(1-\alpha_N)=(\omega+\beta_N)d_{K,N}$.
Thus $r_{N+1}\le d_{K,N}\le1$ and both determinant ratios are
positive, even when one kernel or image is zero dimensional.
The exact original per-step logarithmic allocations are
\begin{equation}\tag{ORE7}
 -\log d_{K,N},\qquad
 -\log r_{N+1}+\log d_{K,N}.
\end{equation}
Their sum is $-\log r_{N+1}$. These formulas compute the
restriction and quotient of the same original metric; they do not
replace the original directional $\alpha_N$ by a rank fraction.

\section{The two boundary responses form one constrained sequence}
Use the original definitions, with the same full $v$,
\[
 U_N=\frac{b_N^*G_NP_Nv}{b_N^*G_Nv},\qquad
 V_N=\frac{b_{N+1}^*G_NP_Nv}{b_{N+1}^*G_Nv}.           \tag{ORE8}
\]
These are exactly NCE14 and EC37; neither row has been exchanged.
Put $s=b^*G_Nv$ and $s_B=b^*G_N(I-P_N)v=(1-V_N)s$.
The update (ORE4) gives $b^*G_{N+1}v=r_{N+1}s$.
The attained quotient inverse update gives
\begin{align}
 Q_{N+1}&=Q_N-
 \frac{Q_N\Lambda b b^*\Lambda^*Q_N}
 {\omega+\beta_N(1-\alpha_N)},\notag\\
 b^*G_{N+1}(I-P_{N+1})v
 &=\frac{r_{N+1}}{d_{K,N}}s_B.                       \tag{ORE9}
\end{align}
At cutoff $N+1$ its first boundary class is literally the same $b$.
Divide the second equality by $r_{N+1}s\ne0$ to obtain
\begin{equation}\tag{ORE10}
 \boxed{\displaystyle
 U_{N+1}=1-\frac{1-V_N}{d_{K,N}}.}
\end{equation}
This is an exact original-source relation between the two response
rows at adjacent cutoffs. It includes all conductor terms through
the actual $H_N$ and every $P_N$. No real part, complex conjugation,
phase or normalization is introduced in this transport.

In particular let $R_N=1-V_N$ and, for $N\ge q$, retain the three
original products at cutoff $N$:
\begin{align}
 P_{B,N}&=\frac{\overline{R_{N-1}}R_N}{d_{K,N-1}},\notag\\
 P_{K,N}&=\left(1-\frac{\overline{R_{N-1}}}{d_{K,N-1}}\right)(1-R_N),
 \qquad P_{\times,N}=1-P_{K,N}-P_{B,N}.                \tag{ORE11}
\end{align}
The first is $(1-\overline U_N)(1-V_N)$ by (ORE10); the second
is $\overline U_NV_N$. Therefore the actual centered pairings are
\begin{equation}\tag{ORE12}
 \Phi_{X,N}=2E_N\Re(\mathfrak c_{N,\lambda}P_{X,N}),\qquad
 \sum_X\Phi_{X,N}=2a\delta E_N,
 \quad X\in\{K,B,\times\}.
\end{equation}
This is the NCE15 current, now with its consecutive response factors
and determinant denominator explicitly linked. The shared phase error
also keeps its signed sum zero by $\sum_XP_{X,N}=1$.
At the first cutoff $U_{q-1}$ remains its original first boundary value;
no fictitious previous quotient of dimension $q$ is inserted.

\section{Exact class energy and observation energy evolution}
Define the original second boundary energy and its actual class loss
\begin{equation}\tag{ORE13}
 p_{2,N}=\frac{|b^*G_Nv|^2}{E_N\beta_N},\qquad
 \eta_N=\xi_Np_{2,N},\qquad
 h_N=\frac{E_{N+1}}{E_N}=1-\eta_N.
\end{equation}
The last equality is direct substitution in (ORE4). Positivity of
$G_{N+1}$ gives $h_N>0$. In fact $0<p_{2,N}\le1$ and
$0<\eta_N<1$. The Cauchy--Schwarz upper bound follows by minimizing
the full positive square of $v-tb$; it uses the original $G_N$.
Applying (ORE9) to the fixed observation value $\Lambda v$ proves
\begin{equation}\tag{ORE14}
 \boxed{(1-\eta_N)(1-\theta_{N+1})
 =1-\theta_N-\frac{\eta_N|1-V_N|^2}{d_{K,N}}.}
\end{equation}
Indeed the quotient energy drop is precisely
$|s_B|^2/[\omega+\beta_N(1-\alpha_N)]$;
divide by $E_N$, use (ORE13) and retain its denominator from (ORE6).
Rearranging gives the same exact update for the kernel component:
\begin{equation}\tag{ORE15}
 \theta_{N+1}-\theta_N
 =\frac{\eta_N}{1-\eta_N}
 \left[\theta_N-1+\frac{|1-V_N|^2}{d_{K,N}}\right].
\end{equation}
Thus a small complete-class increment alone does not control the
change of the original kernel projection: its actual amplification
and its restriction determinant appear in the exact same equation.
The nonnegative quotient energy gives the finite evaluated bound
\begin{equation}\tag{ORE16}
 0\le\frac{\eta_N|1-V_N|^2}{d_{K,N}}\le1-\theta_N.
\end{equation}
If $\Lambda v=0$, both sides of the underlying energy drop vanish,
$\theta_N=1$ and $U_N=V_N=1$ for every cutoff. If $\Lambda v\ne0$,
the quotient energy is strictly positive at every cutoff by $Q_N>0$.
This resolves the possible zero value without deleting that original
class or taking an undefined logarithm.

There is also an exact propagation of the two original boundary energies.
For the unchanged $b=b_{N+1}$, (ORE4) gives
$b^*G_{N+1}b=r_{N+1}\beta_N$ and
$b^*G_{N+1}v=r_{N+1}s$. Hence
\begin{equation}\tag{ORE17}
 \boxed{p_{1,N+1}=
 \frac{r_{N+1}p_{2,N}}{1-(1-r_{N+1})p_{2,N}}.}
\end{equation}
Combining this with the independent original parity calculation
in the companion boundary-energy proof evaluates both energies from
the same consecutive source moments, with no metric switch.

\section{Both windows and the full original coefficient transition}
For any positive sequence $f_N$, direct cancellation of consecutive
ratios proves
\begin{equation}\tag{ORE18}
 \log\frac{f_{2q-1}f_{2q}}{f_{q-1}f_q}
 =\sum_{N=q-1}^{2q-1}w_N\log\frac{f_{N+1}}{f_N},
 \quad w_{q-1}=w_{2q-1}=1,\quad w_N=2\text{ otherwise}.
\end{equation}
Apply it to $\det G_N$, $\det H_N$, $\det Q_N$ and $E_N$.
Their increments are respectively $\log r_{N+1}$,
$\log d_{K,N}$, $\log r_{N+1}-\log d_{K,N}$ and
$\log(1-\eta_N)$. Both original windows and each determinant
frame cancellation are thereby retained. If $\Lambda v\ne0$,
it also applies to $E_N(1-\theta_N)$ with its exact ratio (ORE14).
If $\Lambda v=0$, that energy is identically zero as evaluated above,
and (ORE14), rather than a logarithm, is its receiving equation.

The connection to the original conductor/native chart is explicit.
For the actual retained index set $I_{\rm sel}$, COI11 gives
$L_{I_{\rm sel}}=F_CD_{0I_{\rm sel}}+L_0M_{0I_{\rm sel}}$.
NCRP's complete frame is $\mathcal F=[F_0,F_C,L_{I_{\rm sel}}]$,
an isomorphism onto the full packet. Put $Y=\mathcal F^{-1}I$.
Then $\Lambda\mathcal FY=0$ and the actual kernel metric is exactly
\begin{equation}\tag{ORE19}
 H_N=Y^*\mathcal F^*G_N\mathcal FY.
\end{equation}
Expanding $\mathcal F^*G_N\mathcal F$ retains all nine blocks; expanding
its $L_{I_{\rm sel}}$ block gives all four terms of COI13. Consequently
every ratio $d_{K,N}$ in (ORE10)--(ORE18) uses those full blocks.
The observation in the analytic native coefficient space is
$B_a=Z_a^{\mathsf T}A_{\ell,a}$ of NSO15, with the exact dual metric
square stated there. This ordinary transpose is not changed to a star,
and this native metric is not substituted for $H_N$ in (ORE19).
The relation graph and NSO15 are the connecting maps for the native
eight flags; (ORE19) is their complete arithmetic receiving metric.

\section{A direct quotient-current receiver with no response division}
Put $e_{B,N}=E_N(1-\theta_N)$ and retain the actual positive denominator
and the actual quotient observation numerator
\begin{equation}\tag{ORE20}
 D^B_N=\omega_{N+1}+\beta_N(1-\alpha_N),\qquad
 s^B_N=b_{N+1}^*G_N(I-P_N)v,\qquad
 \Delta^B_N=e_{B,N}-e_{B,N+1}=\frac{|s^B_N|^2}{D^B_N}\ge0.
\end{equation}
The last identity is the undivided version of (ORE14), so it remains
valid even when $s^B_N=0$. At the next cutoff (ORE9) says that the
first quotient boundary numerator is
$\omega_N s^B_{N-1}/D^B_{N-1}$. Substituting this actual numerator
into NCE6 yields, for $N\ge q$,
\begin{equation}\tag{ORE21}
 \boxed{\displaystyle
 \Phi_{B,N}=\frac{2\Re(\overline{s^B_{N-1}}s^B_N)}{D^B_{N-1}}.}
\end{equation}
In particular this computes the exact same original pairing as
(ORE11)--(ORE12), without division by either small full-class
boundary inner product. The class, observation and metric have not
changed. If both drops in (ORE20) are positive, define the actual
unit complex phases $z_N=s^B_N/|s^B_N|$ and put
$c^B_N=\sqrt{D^B_N/D^B_{N-1}}$. Then
\begin{equation}\tag{ORE22}
 \Phi_{B,N}=2c^B_N\sqrt{\Delta^B_{N-1}\Delta^B_N}
                  \Re(\overline z_{N-1}z_N),\qquad
 |\Phi_{B,N}|\le2c^B_N\sqrt{\Delta^B_{N-1}\Delta^B_N}.
\end{equation}
If either drop is zero, (ORE20) makes its numerator zero and (ORE21)
gives $\Phi_{B,N}=0$; no phase is assigned to that zero.
Writing $r^B_{N+1}=\det Q_{N+1}/\det Q_N$, (ORE6) further gives
\begin{equation}\tag{ORE23}
 (c^B_N)^2=\frac{\omega_{N+1}}{\omega_N}
             \frac{r^B_N}{r^B_{N+1}},\qquad
 \sum_{N=L}^{H}\Delta^B_N=e_{B,L}-e_{B,H+1}.
\end{equation}
Thus the coefficient in the undivided current is fixed by the original
arithmetic recurrence and original quotient determinants. For any
finite interval $q\le L\le H$, (ORE22) and $2xy\le x^2+y^2$ give the
explicit finite bound
\[
 \sum_{N=L}^{H}|\Phi_{B,N}|
 \le\left(\max_{L\le N\le H}c^B_N\right)
 \left(\Delta^B_{L-1}+2\sum_{N=L}^{H-1}\Delta^B_N+\Delta^B_H\right)
 \le2\left(\max_{L\le N\le H}c^B_N\right)e_{B,L-1}.
\]
Every coefficient and endpoint is retained. This bound does not assert
a new uniform size for $c^B_N$; (ORE23) is the exact determinant quantity
to calculate. The original individual sign in (ORE22) is precisely the
adjacent observed-numerator phase, rather than the phase of the full
eigenclass alone.

\section{Scope and preceding proofs}
The original quotient, boundary and current formulas used here are
NCE1--21 in the complete \emph{Native class integration} source;
the response energies are JRE1--22. All new equalities above follow
from the full source increment (ORE4) and the displayed attained
minimum, with the fixed original $I,\Lambda,v$ throughout.
They reduce the response sequence and its class-energy evolution to
the same actual restriction determinants used by the allocation
problem. The remaining calculation is the values of those original
directional determinants and the residual response phase. Neither
the evaluated full action AG29 nor an upper bound for $\eta_N$ alone
assigns the separate projected signs.

\begin{thebibliography}{9}
\bibitem{Boyd} Stephen Boyd and Lieven Vandenberghe,
\emph{Convex Optimization}, Cambridge University Press, 2004,
Appendices A.5.5 and C.4.1,
\url{https://web.stanford.edu/~boyd/cvxbook/bv_cvxbook.pdf}.
\bibitem{Hager} W. W. Hager, ``Updating the Inverse of a Matrix'',
\emph{SIAM Review} 31 (1989), 221--239,
\href{https://doi.org/10.1137/1031049}{doi:10.1137/1031049}, equation (2).
\end{thebibliography}
\end{document}
